\label{chap:background}

This chapter provides the scientific and technical foundations required for the thesis.

\section{Catalysis}
\label{sec:catalysis}

Catalysis refers to the acceleration of a chemical reaction by a substance that participates in the reaction pathway but is regenerated during the catalytic cycle. A catalyst therefore changes the kinetic route of a reaction rather than the thermodynamic equilibrium: it can lower the effective activation barrier and allow a reaction to proceed under milder or more technically useful conditions, as schematically shown in Figure~\ref{fig:catalysis-activation-energy}. Catalysts may operate in the same phase as the reactants, referred to as homogeneous catalysis, or in a different phase, referred to as heterogeneous catalysis. In technical chemistry, their value lies not only in increasing reaction rates, but also in directing complex reaction networks toward desired products and maintaining performance over time. Activity, selectivity, and stability are therefore central descriptors of catalyst performance and already show why catalytic systems have to be studied under well-defined reaction and material conditions \cite{fechete2012past,job2021physikalische}.

\begin{figure}[H]
    \centering
    \begin{tikzpicture}[
        >=Latex,
        scale=0.95,
        every node/.style={font=\small}
    ]
        \draw[->, line width=0.8pt] (0,0) -- (8.2,0) node[below] {Reaction coordinate};
        \draw[->, line width=0.8pt] (0,0) -- (0,4.5) node[above, anchor=east] {$E_\mathrm{G}$};

        \draw[gray!35, dashed] (0.9,1.0) -- (7.4,1.0);
        \draw[gray!35, dashed] (6.7,0.25) -- (7.6,0.45);

        \coordinate (reactants) at (0.9,1.0);
        \coordinate (products) at (7.4,0.45);
        \coordinate (bluePeak) at (3.15,3.32);
        \coordinate (redPeak) at (2.55,2.10);

        \draw[blue!70!black, line width=1.3pt]
            (reactants) .. controls (1.7,3.05) and (2.35,3.32) .. (bluePeak)
            .. controls (3.85,3.32) and (4.35,2.25) .. (5.0,1.45)
            .. controls (5.7,0.75) and (6.4,0.45) .. (products);
        \draw[red!75!black, line width=1.3pt]
            (reactants) .. controls (1.55,1.85) and (2.05,2.10) .. (redPeak)
            .. controls (3.25,2.10) and (3.95,1.80) .. (4.7,1.45)
            .. controls (5.5,0.85) and (6.4,0.45) .. (products);

        \draw[<->, blue!70!black, line width=0.8pt] (bluePeak |- reactants) -- (bluePeak);
        \draw[<->, red!75!black, line width=0.8pt] (redPeak |- reactants) -- (redPeak);
        \draw[<->, gray!70!black, line width=0.8pt] (7.75,1.0) -- (7.75,0.45)
            node[midway, right] {$\Delta G$};

        \node[anchor=west] at (0.95,0.72) {Reactants};
        \node[anchor=west] at (6.65,0.23) {Products};
        \node[blue!70!black, anchor=west, fill=white, inner sep=1pt] at (3.32,3.18) {$E_{\mathrm{a}}$ without catalyst};
        \node[red!75!black, anchor=east, fill=white, inner sep=1pt] at (2.42,1.72) {$E_{\mathrm{a}}$ with catalyst};
    \end{tikzpicture}
    \caption{Schematic reaction coordinate diagram illustrating catalyzed and uncatalyzed reaction pathways, adapted from \cite{job2021physikalische}.}
    \label{fig:catalysis-activation-energy}
\end{figure}


\subsection{Catalytic Experiments}
\label{sec:catalytic-experiments}

Catalytic experiments investigate how catalysts are prepared, activated, characterized, and used to transform reactants under defined reaction conditions. The resulting data landscape is therefore not limited to final performance indicators, but spans the full experimental chain from catalyst synthesis to functional analysis \cite{marshall2023digital}.
During the life cycle of a catalyst, a precursor may be transformed into an activated material and subsequently into a spent or deactivated sample. Each stage in this process can yield new characterisation and performance data, which must be linked to the corresponding material state \cite{mendes2025data}.
For this reason, a catalytic result such as a conversion, selectivity, yield, or reaction rate is only interpretable together with information about the catalyst history, reactor configuration, operating conditions, analytical methods, and data processing steps \cite{wulf2021unified}.
Table~\ref{tab:catalytic-data-landscape} summarizes this landscape by grouping representative data and metadata into material, preparation, characterization, reactor, operating, and performance-related categories. These categories are not an exhaustive ontology, but they indicate which contextual links have to be preserved when heterogeneous datasets and kinetic measurements, are interpreted together \cite{mendes2025data,schumann2025nomad,stroemert2026chemdcat}.

\begin{table}[H]
    \centering
    \caption{Representative data and metadata types in catalytic experiments.}
    \label{tab:catalytic-data-landscape}
    \small
    \renewcommand{\arraystretch}{1.2}
    \begin{tabularx}{\textwidth}{p{0.23\textwidth}p{0.35\textwidth}X}
        \toprule
        Experimental area & Typical data and metadata & Role in interpretation \\
        \midrule
        Catalyst material & Composition or formula, metal loading, support material, and unique sample identifier & Identifies the material under study and links results to a specific catalyst sample \cite{schumann2025nomad,salim2026methane}. \\
        Synthesis and treatment & Preparation method, precursor identity, calcination or reduction parameters, and treatment atmosphere & Preserves the catalyst history, since preparation and activation steps define the material state used in testing \cite{mendes2025data,salim2026methane}. \\
        Physicochemical characterization & Surface area, pore properties, particle or crystallite size, and surface composition & Connects structural and chemical descriptors to catalytic behavior and material comparison \cite{schumann2025nomad,salim2026methane}. \\
        Reactor and setup & Reactor type, reactor dimensions, catalyst bed mass, diluent, and reactor filling & Provides the experimental context needed to interpret kinetic measurements and compare setups \cite{wulf2021unified,schumann2025nomad}. \\
        Operating conditions & Reaction temperature and pressure, feed composition, total flow rate, and space velocity or contact time & Defines the process conditions under which performance values were measured \cite{schumann2025nomad,salim2026methane}. \\
        Performance and stability & Conversion, selectivity, yield, reaction rate, and time-on-stream behavior & Reports the functional outcome of the experiment and supports comparison across catalysts and conditions \cite{schumann2025nomad,salim2026methane}. \\
        \bottomrule
    \end{tabularx}
\end{table}


Experimental metadata describe the context in which primary data were generated, processed, and interpreted. In catalytic experiments, this includes sample provenance, instrument settings, calibration procedures, and processing choices. Such information is not just a secondary annotation layer, but rather an integral part of the scientific evidence, because the same numerical activity value can have a different meaning when the catalyst pretreatment or analytical calibration changes \cite{wulf2021unified}. Metadata therefore turn isolated measurements into documented experimental records and provide the basis for reproducing experiments or comparing evaluation studies across different catalysts, reactors, and laboratories \cite{wulf2021unified,salim2026methane}. In automated catalyst testing, they also keep raw data, evaluated data, standard operating procedures, and database entries connected throughout the experimental workflow \cite{moshantaf2024advancing}.

TODO: Give a real example of metadata entries from an experiment and how they affect the interpretation of a result.

\section{Semantic Web Technologies}
\label{sec:semantic-web-technologies}
The previous section showed that catalytic data only become scientifically interpretable when measurements are connected to their experimental context. Semantic Web technologies provide the technical foundation for representing this context in a machine-actionable form. Their central idea is to move from a web of human-readable documents to a web of data whose meaning is explicitly described for computational processing \cite{jannidis2017digitalhumanities}.
In contrast to the free-text documentation of metadata attached to files, semantic approaches explicitly describe entities, relations, and constraints. This transforms data from mere stored information into interconnected statements that can be searched, queried, validated, combined, and reused across repositories and workflows.

\subsection{Ontologies and Controlled Vocabularies}
\label{sec:ontologies-controlled-vocabularies}

An ontology is a formal specification of the concepts and relations that are relevant in a domain. In contrast to an informal glossary, an ontology not only lists terms, but also defines how these terms may be related, for example through class hierarchies, properties, restrictions, and equivalence relations. Ontologies therefore provide a shared conceptual model that can be used by both humans and machines to interpret data consistently \cite{stuckenschmidt2011ontologien}. They can be distinguished by scope: upper ontologies describe very general categories, domain ontologies capture the terminology of a field, and application ontologies adapt these concepts to a concrete tool, workflow, or data model \cite{guarino1998formal,nfdi4cat2024whitepaper}.

The central modelling elements used in this work are classes, individuals, and properties. Classes represent types of entities such as catalysts, reactions, measurements, units, or analytical methods. Individuals represent concrete instances, for example a particular catalyst sample, experiment, extracted value, or dataset record. Properties connect these entities to one another or to literal values such as numbers or labels. Axioms and restrictions can further state subclass relations, allowed value types, or required relations. In knowledge graph construction, the ontology therefore provides the semantic framework: nodes and edges instantiate ontology classes and properties \cite{stuckenschmidt2011ontologien,arp2015building,diaz2025knowledgegraphs}.

The Resource Description Framework (RDF) provides the graph data model for representing such grounded metadata in a knowledge graph as subject--predicate--object triples, while RDF Schema (RDFS) and Web Ontology Language (OWL) add schema-level semantics such as class membership, subclass relations, and logical restrictions \cite{w3c2014rdf11concepts,w3c2014rdfschema,w3c2012owl2primer}.
This distinction is important because information in triple form is not yet reusable metadata by itself. It becomes semantically useful when it is linked to a vocabulary term, class, property, or individual with a stable identifier. Figure~\ref{fig:ontology-grounding-layers} illustrates this schema-instance distinction for catalysis metadata.

\begin{figure}[!htbp]
    \centering
    \resizebox{0.9\textwidth}{!}{%
    \begin{tikzpicture}[
        font=\normalfont,
        >=Stealth,
        classnode/.style={
            ellipse,
            draw=blue!70!black,
            fill=blue!5,
            very thick,
            minimum width=2.65cm,
            minimum height=1.0cm,
            align=center,
            font=\normalfont\small
        },
        instancenode/.style={
            ellipse,
            draw=green!45!black,
            fill=green!9,
            very thick,
            minimum width=2.65cm,
            minimum height=1.0cm,
            align=center,
            font=\normalfont\small
        },
        literalnode/.style={
            rounded rectangle,
            draw=green!45!black,
            fill=green!6,
            very thick,
            minimum width=2.5cm,
            minimum height=0.8cm,
            align=center,
            font=\normalfont\small
        },
        layerbox/.style={
            rounded corners=8pt,
            draw=black!30,
            fill=black!2,
            line width=0.8pt
        },
        schemaedge/.style={->, line width=1.0pt, draw=blue!70!black},
        instanceedge/.style={->, line width=1.0pt, draw=green!45!black},
        typeedge/.style={->, dotted, line width=1.0pt, draw=black!45},
        edgelabel/.style={
            fill=white,
            inner sep=1.6pt,
            font=\normalfont\scriptsize,
            text=black
        },
        typelabel/.style={
            fill=white,
            inner sep=2.0pt,
            font=\normalfont\scriptsize,
            text=black!75
        },
        layerlabel/.style={
            font=\normalfont\Large\bfseries,
            text=black,
            anchor=west
        }
    ]

    \begin{scope}[on background layer]
        \node[layerbox, minimum width=17.2cm, minimum height=4.6cm] at (0,2.65) {};
        \node[layerbox, minimum width=17.2cm, minimum height=4.7cm] at (0,-3.15) {};
    \end{scope}

    \node[layerlabel] at (-8.35,5.25) {Ontology schema layer};
    \node[layerlabel] at (-8.35,-0.45) {Instance layer};

    \node[classnode] (dataset) at (-6.75,3.45) {Measurement\\Dataset};
    \node[classnode] (reaction) at (-1.2,3.55) {Reaction};
    \node[classnode] (catreaction) at (-1.2,1.55) {Catalytic\\Reaction};
    \node[classnode] (catalyst) at (4.0,3.9) {Catalyst};
    \node[classnode] (condition) at (3.65,2.0) {Reaction\\Condition};
    \node[classnode] (substance) at (7.35,2.45) {Substance};

    \draw[schemaedge] (dataset) -- (reaction)
        node[edgelabel, midway, above=2pt] {records result for};
    \draw[schemaedge] (catreaction) -- (reaction);
    \node[edgelabel] at (-2.05,2.5) {subclass of};
    \draw[schemaedge] (reaction) -- (catalyst)
        node[edgelabel, midway, above=3pt] {has catalyst};
    \draw[schemaedge] (reaction) -- (condition);
    \node[edgelabel] at (1.05,2.55) {requires condition};
    \draw[schemaedge] (reaction) to[out=-8,in=165] (substance);
    \node[edgelabel] at (5.9,3.0) {has reactant};
    \draw[schemaedge] (catreaction) to[out=-8,in=202] (substance);
    \node[edgelabel] at (4.75,1.42) {has product};

    \node[instancenode] (dataset1) at (-6.75,-2.25) {Dataset D1};
    \node[instancenode] (r1) at (-1.2,-2.25) {Reaction R1};
    \node[instancenode] (cat1) at (3.65,-1.25) {Iron catalyst};
    \node[instancenode] (co2) at (7.35,-2.05) {CO$_2$};
    \node[instancenode] (methane) at (7.35,-3.35) {CH$_4$};
    \node[literalnode] (conversion) at (-6.75,-4.15) {58\% conversion};
    \node[literalnode] (temp) at (3.65,-4.35) {320 $^\circ$C};

    \draw[instanceedge] (dataset1) -- (r1)
        node[edgelabel, midway, above=2pt] {describes reaction};
    \draw[instanceedge] (dataset1) -- (conversion)
        node[edgelabel, midway, left=2pt] {reports value};
    \draw[instanceedge] (r1) -- (cat1);
    \node[edgelabel] at (1.1,-1.65) {has catalyst};
    \draw[instanceedge] (r1) -- (co2);
    \node[edgelabel] at (4.75,-2.42) {has reactant};
    \draw[instanceedge] (r1) -- (methane);
    \node[edgelabel] at (4.85,-3.0) {has product};
    \draw[instanceedge] (r1) -- (temp)
        node[edgelabel, pos=0.5, below left=1pt] {has temperature};

    \draw[typeedge] (dataset1) -- (dataset);
    \draw[typeedge] (r1) -- (reaction);
    \draw[typeedge] (cat1) -- (catalyst);
    \draw[typeedge] (co2) -- (substance);
    \draw[typeedge] (methane) -- (substance);
    \draw[typeedge] (temp) -- (condition);

    \draw[typeedge] (4.85,-5.05) -- (5.6,-5.05);
    \node[typelabel, rounded rectangle, draw=black!20, anchor=west] at (5.7,-5.05) {rdf:type / instance of};

    \end{tikzpicture}%
    }
    \caption{Schema and instance layers of a catalysis knowledge graph. The ontology schema defines reusable classes and relations, while concrete metadata instances populate the graph as labeled subject--predicate--object relations \cite{stuckenschmidt2011ontologien,diaz2025knowledgegraphs,nfdi4cat2024whitepaper}.}
    \label{fig:ontology-grounding-layers}
\end{figure}



\section{Large Language Models}
\label{sec:large-language-models}
Large Language Models (LLMs) are neural language models that learn statistical regularities from large text corpora and use these learned representations to predict, generate, transform, and classify text \cite{zhao2026surveyllms,brown2020language}. They are part of the broader development of generative artificial intelligence, in which models do not only assign labels to predefined inputs, but generate new sequences such as natural-language explanations, code, tables, or structured records \cite{zhao2026surveyllms,dagdelen2024structured}. The following sections introduce the technical foundations that are most relevant for understanding LLM-based systems: transformer architectures, retrieval-augmented generation, and methods for structured output generation.

\subsection{Transformer Architecture}
\label{subsec:transformer-architecture}

Most current LLMs are based on the transformer architecture, which replaced recurrent sequence processing with attention-based computation \cite{vaswani2017attention,zhao2026surveyllms}. Text is first divided into tokens, which are mapped to numerical vector representations. These token embeddings do not store words as discrete symbols alone, but place them in a continuous vector space in which semantic and syntactic regularities can be represented numerically \cite{devlin2019bert}. Positional information is added so that the model can distinguish the order of tokens, because the attention mechanism itself processes the sequence as a set of vector representations \cite{vaswani2017attention}.

The central mechanism of a transformer is self-attention \cite{vaswani2017attention}. For each token, the model computes query, key, and value vectors. Attention weights are obtained by comparing queries with keys, and these weights determine how strongly the value vectors of other tokens contribute to the updated representation of the current token. Multi-head attention repeats this operation in several learned subspaces, allowing the model to represent different types of relationships in parallel. Feed-forward layers, residual connections, layer normalisation, and stacked transformer blocks then transform these contextualised token representations into deeper representations that can support language understanding and generation \cite{vaswani2017attention,devlin2019bert}.

\begin{figure}[H]
    \centering
    \resizebox{\textwidth}{!}{%
    \begin{tikzpicture}[
        font=\normalfont,
        >=Stealth,
        token/.style={
            rounded rectangle,
            draw=black!55,
            fill=black!4,
            line width=0.8pt,
            minimum width=1.75cm,
            minimum height=0.75cm,
            align=center,
            font=\normalfont\small
        },
        vector/.style={
            rounded rectangle,
            draw=blue!70!black,
            fill=blue!6,
            line width=0.9pt,
            minimum width=2.15cm,
            minimum height=0.75cm,
            align=center,
            font=\normalfont\small
        },
        block/.style={
            rounded rectangle,
            draw=green!45!black,
            fill=green!7,
            line width=1.0pt,
            minimum width=3.0cm,
            minimum height=0.95cm,
            align=center,
            font=\normalfont\small
        },
        outblock/.style={
            rounded rectangle,
            draw=orange!80!black,
            fill=orange!9,
            line width=1.0pt,
            minimum width=2.6cm,
            minimum height=0.85cm,
            align=center,
            font=\normalfont\small
        },
        arrow/.style={->, line width=0.9pt, draw=black!65},
        attn/.style={->, line width=0.9pt, draw=red!70!black, bend left=28}
    ]

    \node[token] (tok1) at (-6.3,3.0) {Token 1};
    \node[token] (tok2) at (-4.3,3.0) {Token 2};
    \node[token] (tok3) at (-2.3,3.0) {Token 3};
    \node[token] (tok4) at (-0.3,3.0) {Token 4};

    \node[vector] (emb1) at (-6.3,1.55) {Embedding\\+ position};
    \node[vector] (emb2) at (-4.3,1.55) {Embedding\\+ position};
    \node[vector] (emb3) at (-2.3,1.55) {Embedding\\+ position};
    \node[vector] (emb4) at (-0.3,1.55) {Embedding\\+ position};

    \node[block] (qkv) at (3.0,1.55) {Query, key, value\\projections};
    \node[block] (attention) at (3.0,-0.1) {Multi-head\\self-attention};
    \node[block] (ffn) at (3.0,-1.75) {Feed-forward layer\\+ residual paths};
    \node[outblock] (next) at (6.8,-1.75) {Next-token\\probabilities};

    \draw[arrow] (tok1) -- (emb1);
    \draw[arrow] (tok2) -- (emb2);
    \draw[arrow] (tok3) -- (emb3);
    \draw[arrow] (tok4) -- (emb4);

    \draw[arrow] (emb1) -- (qkv);
    \draw[arrow] (emb2) -- (qkv);
    \draw[arrow] (emb3) -- (qkv);
    \draw[arrow] (emb4) -- (qkv);
    \draw[arrow] (qkv) -- (attention);
    \draw[arrow] (attention) -- (ffn);
    \draw[arrow] (ffn) -- (next);

    \draw[attn] (tok1.north) to (tok3.north);
    \draw[attn] (tok2.north) to (tok4.north);
    \draw[attn] (tok4.north) to[bend left=42] (tok1.north);

    \end{tikzpicture}%
    }
    \caption{Placeholder sketch of a transformer block and self-attention mechanism. Tokens are mapped to embeddings with positional information, transformed into query, key, and value vectors, and contextualised through multi-head self-attention before the decoder produces next-token probabilities \cite{vaswani2017attention}.}
    \label{fig:transformer-attention-placeholder}
\end{figure}


Transformer models are commonly grouped into encoder-only, decoder-only, and encoder--decoder architectures \cite{zhao2026surveyllms}. Encoder-only models, such as BERT-like systems, process the input bidirectionally and are well suited for representation learning, classification, and span-level tasks \cite{devlin2019bert,zhao2026surveyllms}. Decoder-only models, such as GPT-style systems, generate text autoregressively by predicting the next token from the preceding context and then repeating this step until a complete output has been produced \cite{brown2020language,zhao2026surveyllms}. Encoder--decoder models combine an input encoder with an output decoder and are commonly used for sequence-to-sequence transformations \cite{vaswani2017attention,zhao2026surveyllms}.

The generative behaviour of many LLMs is based on language modelling \cite{zhao2026surveyllms,brown2020language}. During pretraining, the model learns from large text collections with a self-supervised objective, such as predicting missing or subsequent tokens. Although this objective is simple, large-scale pretraining can produce models that capture linguistic regularities, factual associations, and task-like patterns from text \cite{brown2020language,zhao2026surveyllms}. A pretrained model, however, does not automatically behave like a helpful instruction-following assistant \cite{ouyang2022training}. Later training stages therefore adapt the model to user-facing tasks. Instruction tuning trains the model on examples of instructions and desired responses, while reinforcement learning from human feedback can further align outputs with human preferences and task intent \cite{ouyang2022training,zhao2026surveyllms}.

In-context learning describes the ability of large language models to adapt their behaviour from information provided directly in the prompt, without updating the model weights \cite{brown2020language,zhao2026surveyllms}. A prompt may contain a task description, formatting rules, background information, or a small number of examples. The model processes these elements together with the current input as one context and generates a continuation that follows the patterns implied by that context \cite{brown2020language,liu2024lostmiddle}. This makes the prompt an operational part of the model interaction: it can specify what kind of output is expected, how ambiguities should be handled, and which examples should guide the response, while the underlying model parameters remain unchanged.

Text embeddings are also important beyond the internal token representations of a transformer \cite{reimers2019sentencebert}. Encoder models and specialised sentence-transformer models can map words, sentences, paragraphs, or documents into dense vectors that capture semantic similarity \cite{devlin2019bert,reimers2019sentencebert}. Such embeddings make it possible to compare texts by vector distance, cluster related passages, or retrieve semantically similar documents \cite{reimers2019sentencebert,dreger2025knowledgegraph}. They therefore provide a bridge between symbolic text and numerical search procedures, which is one reason why embeddings are frequently used together with retrieval-based LLM systems \cite{gao2023ragsurvey}.

\subsection{Retrieval Augmented Generation}
\label{subsec:retrieval-augmented-generation}

Retrieval-Augmented Generation (RAG) extends a generative language model with an explicit retrieval step \cite{lewis2020retrieval,gao2023ragsurvey}. Instead of relying only on information stored implicitly in the model parameters during training, the system first searches an external collection for relevant documents, passages, database entries, schema descriptions, or examples. The retrieved material is then inserted into the model context, and the model generates an answer conditioned on both the user input and the supplied external evidence \cite{lewis2020retrieval}.

\begin{figure}[H]
    \centering
    \resizebox{\textwidth}{!}{%
    \begin{tikzpicture}[
        font=\normalfont,
        >=Stealth,
        usernode/.style={
            rounded rectangle,
            draw=blue!70!black,
            fill=blue!6,
            line width=0.9pt,
            text width=1.85cm,
            minimum height=0.9cm,
            inner sep=4pt,
            align=center,
            font=\normalfont\small
        },
        processnode/.style={
            rounded rectangle,
            draw=green!45!black,
            fill=green!8,
            line width=1.0pt,
            text width=2.25cm,
            minimum height=0.95cm,
            inner sep=4pt,
            align=center,
            font=\normalfont\small
        },
        contextnode/.style={
            rounded rectangle,
            draw=orange!80!black,
            fill=orange!10,
            line width=1.0pt,
            text width=2.25cm,
            minimum height=0.95cm,
            inner sep=4pt,
            align=center,
            font=\normalfont\small
        },
        storenode/.style={
            cylinder,
            shape border rotate=90,
            aspect=0.25,
            draw=black!55,
            fill=black!5,
            line width=0.9pt,
            text width=2.2cm,
            minimum height=1.25cm,
            inner sep=4pt,
            align=center,
            font=\normalfont\small
        },
        phasebox/.style={
            rounded corners=8pt,
            draw=black!25,
            fill=black!2,
            line width=0.8pt
        },
        phaselabel/.style={
            font=\normalfont\bfseries,
            text=black!75,
            anchor=west
        },
        arrow/.style={->, line width=0.9pt, draw=black!65},
        queryarrow/.style={->, line width=0.9pt, draw=blue!70!black},
        retrievalarrow/.style={->, line width=0.9pt, draw=green!45!black},
        dashedarrow/.style={->, dashed, line width=0.9pt, draw=black!55}
    ]

    \begin{scope}[on background layer]
        \node[phasebox, minimum width=15.3cm, minimum height=1.85cm] at (0,-2.4) {};
        \node[phasebox, minimum width=15.3cm, minimum height=2.25cm] at (0,0.8) {};
    \end{scope}

    \node[phaselabel] at (-7.45,1.85) {Inference time};
    \node[phaselabel] at (-7.45,-1.55) {Indexing time};

    \node[usernode] (query) at (-6.25,0.8) {User\\query};
    \node[processnode] (retriever) at (-3.75,0.8) {Retriever};
    \node[contextnode] (passages) at (-1.1,0.8) {Relevant\\passages};
    \node[processnode] (prompt) at (1.65,0.8) {Prompt:\\query + context};
    \node[processnode] (llm) at (4.4,0.8) {Generative\\LLM};
    \node[contextnode] (answer) at (6.9,0.8) {Grounded\\answer};

    \node[usernode] (docs) at (-6.25,-2.4) {Source\\documents};
    \node[processnode] (chunking) at (-3.75,-2.4) {Chunking and\\preprocessing};
    \node[processnode] (features) at (-1.1,-2.4) {Sparse or dense\\features};
    \node[storenode] (index) at (1.65,-2.4) {External\\knowledge index};

    \draw[queryarrow] (query) -- (retriever);
    \draw[retrievalarrow] (retriever) -- (passages);
    \draw[arrow] (passages) -- (prompt);
    \draw[arrow] (prompt) -- (llm);
    \draw[arrow] (llm) -- (answer);

    \draw[arrow] (docs) -- (chunking);
    \draw[arrow] (chunking) -- (features);
    \draw[arrow] (features) -- (index);

    \draw[retrievalarrow] (index.north) to[out=100,in=-70] (retriever.south);
    \draw[dashedarrow] (query.south east) to[out=-35,in=185] (prompt.west);

    \end{tikzpicture}%
    }
    \caption{Placeholder sketch of a retrieval-augmented generation architecture. Source documents are indexed for search; at inference time, a user query is used to retrieve relevant context, which is combined with the query before generation \cite{lewis2020retrieval}.}
    \label{fig:rag-architecture-placeholder}
\end{figure}


This separation is conceptually important because model parameters and external knowledge serve different roles \cite{lewis2020retrieval,gao2023ragsurvey}. Parameter knowledge is compressed into the weights during training and is difficult to inspect, update, or restrict after deployment. Retrieved context, by contrast, can be selected at inference time, can reflect a curated or current knowledge source, and can be shown together with the generated response. RAG therefore reduces the need for the LLM to answer purely from memory and gives the generation process access to explicit supporting information \cite{gao2023ragsurvey}.

A common RAG pipeline consists of indexing, retrieval, context construction, and generation \cite{gao2023ragsurvey,zhao2026surveyllms}. During indexing, source documents are split into smaller units and represented for search, often by sparse lexical features, dense embeddings, or a combination of both. During retrieval, the system selects the most relevant units for the current query. These units are then placed into a prompt or context window together with the user's request. The LLM finally generates an output that is conditioned on this retrieved material \cite{lewis2020retrieval,gao2023ragsurvey}. The quality of such a system therefore depends not only on the language model, but also on the relevance, granularity, and ordering of the retrieved context \cite{gao2023ragsurvey,liu2024lostmiddle}.

RAG does not remove all limitations of LLMs \cite{gao2023ragsurvey,asai2024selfrag}. A model may ignore relevant retrieved evidence, overinterpret ambiguous passages, or generate unsupported details if the prompt and validation procedure do not constrain the response sufficiently \cite{liu2024lostmiddle,asai2024selfrag}. Nevertheless, RAG changes the basic information source available to the model: the answer no longer has to be generated only from the statistical associations stored in the parameters, but can be grounded in information supplied at inference time \cite{lewis2020retrieval,gao2023ragsurvey}.

\subsection{Structured Output}
\label{subsec:structured-output}

LLMs generate text token by token, but many downstream workflows require outputs that can be processed as structured data. Structured-output techniques address this by constraining the model interaction before, during, or after generation \cite{schilling2025text,dagdelen2024structured}.

The simplest approach is to specify the desired format directly in the prompt, for example JSON, YAML, XML, Markdown tables, or subject--predicate--object triples \cite{dagdelen2024structured}. The prompt may define fields, allowed values, nesting rules, or examples. However, this only guides the model: the structure is still generated as ordinary text and may therefore be malformed, incomplete, or inconsistent with the requested schema \cite{dagdelen2024structured,caufield2024spires}.

A stronger interface-level approach is tool or function calling \cite{schick2023toolformer,schilling2025text}. Here, the model is asked to select a predefined tool and fill its named parameters. This makes the boundary between natural-language reasoning and machine-readable action more explicit, because arguments can be checked against the tool specification before execution or acceptance \cite{schick2023toolformer,caufield2024spires}. Nevertheless, valid argument structure does not guarantee that the generated values are factually correct or grounded in the input \cite{caufield2024spires,yao2023react}.

A third approach is output parsing and validation \cite{dreger2025knowledgegraph,schilling2025text}. The generated response is treated as a candidate structure and is accepted only if it can be parsed and validated against constraints such as schemas or SHACL shapes \cite{w3c2017shacl,dreger2025knowledgegraph}. This can reject malformed records, enforce required fields, and normalise datatypes. However, validation mainly checks processability; semantic errors can remain when a syntactically valid output contains unsupported or inappropriate content \cite{caufield2024spires,dreger2025knowledgegraph}.

\begin{table}[H]
    \centering
    \caption{Limitations of structured output approaches.}
    \label{tab:structured-output-approaches}
    \small
    \renewcommand{\arraystretch}{1.2}
    \begin{tabularx}{\textwidth}{
        >{\raggedright\arraybackslash}p{0.28\textwidth}
        >{\raggedright\arraybackslash}X
    }
        \toprule
        Approach & Limitations \\
        \midrule
        Prompt instructions & Formatting is guided but not guaranteed; malformed structures or unsupported content can still occur and therefore have to be checked after generation \cite{dagdelen2024structured,caufield2024spires}. \\
        Tool or function calling & The interface can enforce structure, but it does not by itself ensure that argument values are factually correct or grounded in the input \cite{caufield2024spires,rankovic2025uncertainty}. \\
        Output parsing and validation & Validation can reject malformed records, but semantic errors may remain when the output is syntactically valid yet unsupported or inappropriate \cite{dagdelen2024structured,rankovic2025uncertainty}. \\
        \bottomrule
    \end{tabularx}
\end{table}

Table~\ref{tab:structured-output-approaches} compares these three approaches according to where they constrain the workflow \cite{schilling2025text}. Prompt instructions guide the model before generation, tool or function calling constrains the interaction interface during generation, and parsing or validation checks the result after generation. The table also highlights their shared limitation: structured-output methods make LLM responses more predictable and machine-actionable, but they do not by themselves guarantee completeness, factual correctness, or grounding in the source material \cite{dreger2025knowledgegraph}.
